\newcommand{\tipoRequerimento}{Solicitação de Dados da CIP}

\section{FUNDAMENTAÇÃO LEGAL DA REQUISIÇÃO}

\begin{enumerate}[resume]
\item
  Para embasar a requisição, utiliza-se as seguintes disposições da Resolução Normativa nº 1.000/2021 (REN nº 1.000/2021) da ANEEL:
  \begin{enumerate}
  \item
    Demonstrativos e Memória de Cálculo (Art. 472, § 1º): A distribuidora tem a obrigação de fornecer ao município os demonstrativos detalhados e a memória de cálculo do faturamento da iluminação pública. Prazos de Atendimento (Art. 408):
    \begin{enumerate}
      \item  
      5 dias úteis para respostas que não exigem visita técnica.
      \item 
      10 dias úteis para os demais casos.
    \end{enumerate}
  \item 
  Prazo Geral (Art. 409): Caso não haja um prazo específico na regulamentação da ANEEL para uma solicitação, a distribuidora deve atendê-la em até 30 dias.
  \end{enumerate}
\end{enumerate}

\section{DOS REQUERIMENTOS}

\begin{enumerate}[resume]
\item Cópia do Convênio firmado entre o Município e a distribuidora para operacionalização da arrecadação da CIP.
\item Cópia do instrumento contratual específico ou convênio adicional, caso exista, que trate da cobrança da Taxa ou Contribuição de Iluminação Pública pela distribuidora em nome do Município.
\item Cópia da Lei Municipal que instituiu a Contribuição para Custeio da Iluminação Pública (CIP ou COSIP) e alterações legislativas em vigor.
\item Demonstrativo mensal da Arrecadação e Aplicação da Contribuição de Iluminação Pública (CIP), emitido regularmente pela distribuidora e encaminhado ao Município, contendo:
\begin{enumerate}[resume]
  \item Valor total arrecadado;
  \item Valor total faturado da iluminação pública;
  \item Valor efetivamente aplicado;
  \item Valor repassado ao Município.
\end{enumerate}

\item Relatório mensal contendo a relação de clientes e os valores faturados a título de CIP, abrangendo os últimos cinco anos, com, no mínimo, as seguintes informações:
\begin{enumerate}[resume]
  \item Valor cobrado a título de CIP;
  \item Base de cálculo aplicada;
  \item Alíquota incidente;
  \item Data do faturamento (mês de referência);
  \item Número da fatura correspondente;
  \item Classificação tarifária (classe e subclasse);
  \item Consumo mensal em kWh.
\end{enumerate}

\item Relatório mensal contendo a relação de clientes e os valores efetivamente arrecadados a título de CIP, também dos últimos cinco anos, incluindo no mínimo:
\begin{enumerate}[resume]
  \item Mês de referência do faturamento associado;
  \item Número da fatura correspondente;
  \item Número da unidade consumidora (UC);
  \item Valor arrecadado;
  \item Data da arrecadação.
\end{enumerate}

\item Relatório mensal com a identificação de todos os contribuintes inadimplentes em relação à CIP, referente aos últimos cinco anos, devendo conter, no mínimo, as seguintes informações:
\begin{enumerate}[resume]
  \item Mês de referência da fatura;
  \item Número da fatura;
  \item Número da unidade consumidora (UC);
  \item Valor da CIP faturado;
  \item Consumo em kWh.
\end{enumerate}

\item Solicitamos que os relatórios solicitados nos itens 9 a 11 referente a relação de clientes e valores efetivamente arrecadados, inadimplentes e faturados a título da Contribuição de Iluminação Pública devem ser disponibilizados em formato editável, preferencialmente .xlsx, .csv ou .txt, a fim de viabilizar a análise técnica e cruzamento de dados.
\end{enumerate}

